% =====================================================================
%  Appendix section: the deployed forecast-driven target controller.
%  Body content only; \input from the pseudocode appendix chapter so it
%  numbers under that chapter. The abridged decide() lives at
%  \cref{bm:lst-autoscaler}; this section gives the full target rule.
% =====================================================================

\section{Forecast-driven target controller}\label{bm:alg:autoscaler}

The autoscaler ships two decision rules behind one dispatch. The original rule,
abridged at \cref{bm:lst-autoscaler}, is a single-step ``bang-bang'' controller:
it moves the pool by exactly one backend per action and only when the predicted
load crosses the current capacity. With a provisioning warm-up and a cooldown
timer, that caps the pool's slew rate at one instance per cooldown window, which
is why a sharp arrival burst that needs several instances at once cannot be
served in time. The deployed rule is the \texttt{target} controller shown below.
It sizes the pool to a capacity-based target directly and may jump several
instances in a single action, removing the per-action slew cap.

The evidence that this controller, rather than the forecast, was the binding fix
comes from the strategy benchmark. The shipped moving-average forecaster sets the
predicted rate to the trailing mean and carries no forward lead, so the
MA-predictive strategy and the purely reactive strategy score identically
($S2 \equiv S3$ to the digit in every profile). Holding that same signal fixed
and swapping the single-step rule for the \texttt{target} controller lifts
attainment from \textbf{77.2\%} to \textbf{98.3\%}. The lever was the control
law, not a sharper prediction.

\subsection*{Sizing law and clamp}

The controller sizes the pool from a capacity model. Writing $\hat{r}$ for the
predicted request rate, $c$ for the per-instance capacity in requests per second,
$h$ for the fractional headroom, and $[N_{\min}, N_{\max}]$ for the policy
bounds, the target backend count is

\begin{equation}\label{bm:eq:target}
  N^{\star} = \operatorname{clamp}\!\left(
    \left\lceil \frac{\hat{r}\,(1+h)}{c} \right\rceil,
    \; N_{\min}, \; N_{\max}
  \right),
  \qquad
  \operatorname{clamp}(x, l, u) = \max\!\big(l, \min(u, x)\big).
\end{equation}

The ceiling and the clamp do the structural work: the ceiling rounds a
fractional demand up to whole instances, and the clamp pins the result inside the
host budget. The upper bound $N_{\max}$ is the real ceiling, not a soft target:
in this deployment $N_{\max}=5$ because an eight-backend capacity probe on the
same host ran out of memory, so the budget is a hardware fact rather than a
tuning preference.

\subsection*{Direction, dead-band, and asymmetric cooldowns}

The controller is asymmetric by direction. Scale-out is sized on a robust
offered-rate signal $\hat{r}_{\text{out}} = \max(\hat{r}, \hat{r}_{\text{band}})$,
where $\hat{r}_{\text{band}}$ is the forecast upper band, and scale-in stays on
the served point estimate $\hat{r}$. The upper band is an offered-rate proxy, not
a measurement of true arrivals: in a closed loop the served rate collapses while
the pool is shedding, so sizing growth on the served estimate alone would let the
loop drain instead of grow. Biasing scale-out toward the wider band breaks that
shed-feedback trap; sizing scale-in on the point estimate keeps the drain
conservative.

A decision fires only after two further gates clear. The dead-band suppresses
flapping around an instance boundary: the pool sheds one backend only if the
smaller pool would still cover demand with the full headroom plus an extra slack
band, that is, when

\begin{equation}\label{bm:eq:deadband}
  (N-1)\,c \;\ge\; \hat{r}\,(1 + h + d),
\end{equation}

with $N$ the current count and $d$ the scale-in dead-band. Each direction then
passes an independent cooldown clock. The two clocks are tracked separately so a
recent scale-in never blocks an urgent scale-out, which is what makes the ``fast
out, slow in'' posture possible: the scale-out cooldown defaults to zero so a
spike is answered immediately, while the scale-in cooldown is long so the pool
drains slowly and does not give back the capacity it just earned. Formally, with
$\Delta_{\text{out}}$ and $\Delta_{\text{in}}$ the seconds since the last action
of each direction and $\theta_{\text{out}} \ll \theta_{\text{in}}$ the two
cooldown windows, an otherwise-warranted action is gated by

\begin{equation}\label{bm:eq:cooldown}
  \text{scale\_out permitted} \iff \Delta_{\text{out}} \ge \theta_{\text{out}},
  \qquad
  \text{scale\_in permitted} \iff \Delta_{\text{in}} \ge \theta_{\text{in}}.
\end{equation}

\Cref{bm:lst:target} gives the full controller: the capacity sizing of both
directions, the clamp, the dead-band, the asymmetric cooldown gates, and the
bounded multi-step jump.

\begin{lstlisting}[language=Python,caption={Forecast-driven \texttt{target} controller: capacity sizing, clamp, dead-band, and asymmetric cooldowns. The actuation loop that drives the pool toward the returned count lives in the run loop.},label={bm:lst:target}]
def target_for_load(load_rps, policy):
    # Instances needed to serve load_rps, clamped to [min, max].
    c = policy.per_instance_capacity_rps
    if c <= 0:
        return clamp(policy.min_backends, policy.min_backends, policy.max_backends)
    eff = max(0.0, load_rps) * (1.0 + policy.headroom)
    need = ceil(eff / c)
    return clamp(need, policy.min_backends, policy.max_backends)


def decide_target(predicted_rps, current_count, policy,
                  seconds_since_scale_out, seconds_since_scale_in,
                  offered_rps=None):
    c = policy.per_instance_capacity_rps
    if c <= 0:
        return NoOp(current_count)  # refuse to scale on an invalid capacity

    # Scale-out sizes on the robust upper-band signal so shedding cannot
    # starve the loop; scale-in keeps sizing on the served point estimate.
    out_rps    = predicted_rps if offered_rps is None else max(predicted_rps, offered_rps)
    out_target = target_for_load(out_rps, policy)
    target     = target_for_load(predicted_rps, policy)

    # -- scale OUT: grow toward out_target, fast cooldown ----------------
    if out_target > current_count:
        if (seconds_since_scale_out is not None
                and seconds_since_scale_out < policy.scale_out_cooldown_s):
            return NoOp(current_count)  # scale-out cooldown active
        step = out_target - current_count
        if policy.max_step_out > 0:
            step = min(step, policy.max_step_out)   # bounded multi-step jump
        new = clamp(current_count + step, policy.min_backends, policy.max_backends)
        return ScaleOut(new)

    # -- scale IN: shed cautiously, slow cooldown, with dead-band -------
    if target < current_count:
        shed_floor = predicted_rps * (1.0 + policy.headroom + policy.scale_in_deadband)
        if (current_count - 1) * c < shed_floor:
            return NoOp(current_count)  # dead-band: shedding would breach slack
        if (seconds_since_scale_in is not None
                and seconds_since_scale_in < policy.scale_in_cooldown_s):
            return NoOp(current_count)  # scale-in cooldown active
        step = current_count - target
        if policy.max_step_in > 0:
            step = min(step, policy.max_step_in)    # drain one at a time
        new = clamp(current_count - step, policy.min_backends, policy.max_backends)
        return ScaleIn(new)

    # -- hold: target matches the pool (dead-band) ----------------------
    return NoOp(current_count)
\end{lstlisting}

\Cref{bm:tab:target-params} lists the controller's parameters with the deployed
defaults. The bounds and capacity are read from the live policy, so a runtime
policy reload moves them on the next decision; the sizing, cooldown, step, and
dead-band knobs are deploy-time settings.

\begin{table}[ht]
  \centering
  \caption[Target controller parameters and defaults]{\texttt{target} controller parameters and deployed defaults.}
  \label{bm:tab:target-params}
  \small
  \begin{tabularx}{\linewidth}{@{}l l Y@{}}
    \toprule
    \textbf{Parameter} & \textbf{Default} & \textbf{Role} \\
    \midrule
    \texttt{per\_instance\_capacity\_rps} & $100$ & Per-backend service capacity $c$ in the sizing law. \\
    \addlinespace
    \texttt{min\_backends} & $1$ & Lower clamp bound $N_{\min}$. \\
    \texttt{max\_backends} & $5$ & Upper clamp bound $N_{\max}$, the host budget. \\
    \addlinespace
    \texttt{headroom} & $0.15$ & Fractional safety margin $h$ on predicted load. \\
    \texttt{scale\_in\_deadband} & $0.15$ & Extra slack $d$ required before shedding a backend. \\
    \addlinespace
    \texttt{scale\_out\_cooldown\_s} & $0$ & Minimum seconds between scale-out actions (fast out). \\
    \texttt{scale\_in\_cooldown\_s} & $120$ & Minimum seconds between scale-in actions (slow in). \\
    \addlinespace
    \texttt{max\_step\_out} & $0$ & Cap on instances added per action ($0$ = jump straight to target). \\
    \texttt{max\_step\_in} & $1$ & Cap on instances removed per action (drain one at a time). \\
    \bottomrule
  \end{tabularx}
\end{table}

The capacity, $N_{\min}$, and $N_{\max}$ defaults are read from
\texttt{config/policy.yaml}; the remaining knobs mirror the controller's own
defaults and are set at deploy time.
